\subsection{A Joint State-Space Model for Perceived Utility, Its Mediators, and Weekly Emissions}
\label{sec:joint_model}

We model the evolution of latent perceived utility jointly with the observables
that it generates: the within-week engagement mediators (unprompted page views
$M_{w,d,t}^{E,PV}$, Fitbit wear $M_{w,d}^{E,FW}$, and daily check-in completion
$M_{w,d}^{E,PJ}$), weekly check-in completion $J_w^{\mathrm{week}}$, and the two
perceived-utility survey items $U_{w,1}$ and $U_{w,2}$. The model is fit
independently for each of the $N = 35$ participants; no parameter is shared
across participants, so we suppress the participant index throughout, and every
coefficient, latent state, and residual variance below is understood to be
participant-specific. Weeks are indexed $w = 1, \ldots, W$ ($W = 12$), days
$d = 1, \ldots, 7$, and decision times $t \in \{1, 2\}$.

The model factorizes as a conditionally specified state-space model.
Index $w$ labels the week that is about to unfold: $E_w$, $Y_w$,
$J_w^{\mathrm{week}}$, $U_{w,1}$, and $U_{w,2}$ sit at the end of week
$w-1$ / start of week $w$ (in the MRT table, $J_w^{\mathrm{week}}$ is
\texttt{week_present_lastweek}). Then $J_w^{\mathrm{week}} \mid E_w$ and
the within-week mediators $M_w \mid E_w, J_w^{\mathrm{week}}$ are generated,
and at the end of week $w$ the transition writes $E_{w+1}$, $Y_{w+1}$,
$J_{w+1}^{\mathrm{week}}$. The initial state is fixed at $E_1 = 2$ for all
participants, anchoring the location and scale of the latent process at study
entry.

\paragraph{Latent transition.}
For $w = 2, \ldots, W$,
\begin{align}
E_w
=
g^{E}(S_{w-1}^{E})^\top \alpha^{E}
+
\epsilon_w^E,
\qquad
\epsilon_w^E
\sim
\mathcal{N}(0, \sigma_E^2),
\label{eq:ssm-ew}
\end{align}
where $g^{E}(S_{w-1}^{E})$ stacks an intercept, the previous state $E_{w-1}$,
and the realized weekly mediator summaries $PV_{w-1}$, $FW_{w-1}$, and
$PJ_{w-1}$. Writing $\alpha^{E} = (a_0, a_1, a_2, a_3, a_4)^\top$, the
coefficient $a_1$ captures week-to-week persistence, and $a_2, a_3, a_4$
capture the feedback of realized engagement on the following week's perceived
utility. The summaries use fixed denominators,
\begin{align}
PV_w
=
\tfrac{1}{14}
\sum_{d=1}^{7} \sum_{t=1}^{2} M_{w,d,t}^{E,PV},
\qquad
FW_w
=
\tfrac{1}{7}
\sum_{d=1}^{7} M_{w,d}^{E,FW},
\qquad
PJ_w
=
\tfrac{1}{7}
\sum_{d=1}^{7} M_{w,d}^{E,PJ},
\label{eq:pu-weekly-summaries}
\end{align}
with missing slots contributing zero: a missing wear flag is treated as
non-wear and a missing check-in as non-completion, so that the summaries
coincide exactly with the quantities the simulator regenerates in closed loop.
For a week with no observed data, the summaries entering \eqref{eq:ssm-ew} are
imputed by the participant's own mean; such weeks contribute no emission terms
to the likelihood.

\paragraph{Weekly emissions.}
For $w = 1, \ldots, W$, the opening check-in is drawn from the arriving state,
\begin{align}
J_w^{\mathrm{week}} \mid E_w
&\sim
\operatorname{Bernoulli}\!\bigl(\pi_w^{J}\bigr),
\qquad
\operatorname{logit}\bigl(\pi_w^{J}\bigr) = b_0 + b_1 E_w,
\nonumber\\
U_{w,j} \mid E_w,\, J_w^{\mathrm{week}} = 1
&\sim
\mathcal{N}\!\bigl(c_{0,j} + c_{1,j} E_w,\; \sigma_{U_j}^2\bigr),
\qquad j \in \{1, 2\}.
\label{eq:pu-weekly-emissions}
\end{align}
The survey items are administered within the weekly check-in and are therefore
observed if and only if $J_w^{\mathrm{week}} = 1$; because the completion
process is itself modeled as a function of $E_w$, this missingness mechanism is
fully accounted for by the joint likelihood.
The last Sunday of week $W$ is $J_{W+1}^{\mathrm{week}}$: it is scored on the
terminal predictive of $E_{W+1}$ after the final Gaussian transition, with no
thirteenth mediator week.

\paragraph{Opening check-in on the engagement mediators.}
The same $J_w^{\mathrm{week}}$ that is drawn from $E_w$ may shift that week's
engagement. Each mediator linear predictor below therefore includes
\begin{align}
q^{k}\bigl(E_w, S\bigr)
=
J_{w}^{\mathrm{week}}
\bigl(
q_0^{k} + q_E^{k}\, E_w + q_B^{k}\, B
\bigr),
\label{eq:pu-jw}
\end{align}
where $B$ denotes recent burden and $k$ indexes the mediator model. A missing
opening $J_w^{\mathrm{week}}$ is omitted from the Bernoulli likelihood and
coded as 0 on the mediator query (including week 1).

\paragraph{Decision-time page views (two-part hurdle).}
Let $O_{w,d,t} = \mathbb{1}\{\text{page-view count} > 0\}$ denote occurrence,
and let $M_{w,d,t}^{E,PV}$ denote the standardized log page-view value, with
zero-count slots placed at a fixed floor below the smallest positive value.
Occurrence follows a logistic model and, conditional on occurrence, intensity
follows a Gaussian model on the standardized log scale:
\begin{align}
O_{w,d,t} \mid E_w,\, J_w^{\mathrm{week}}
&\sim
\operatorname{Bernoulli}\!\bigl(\pi_{w,d,t}^{O}\bigr),
&
\operatorname{logit}\bigl(\pi_{w,d,t}^{O}\bigr)
&=
\eta_{w,d,t}^{O},
\nonumber\\
M_{w,d,t}^{E,PV} \mid E_w,\, J_w^{\mathrm{week}},\, O_{w,d,t} = 1
&\sim
\mathcal{N}\!\bigl(\eta_{w,d,t}^{I},\, \sigma_{PV}^2\bigr),
\label{eq:pu-pv-model}
\end{align}
where the two linear predictors share a common functional form with separate
coefficient vectors,
\begin{align}
\eta_{w,d,t}^{k}
=
g^{k}(S_{w,d,t})^\top \alpha^{k}
+
A_{w,d,t}\, f^{k}(S_{w,d,t})^\top \beta^{k}
+
q^{k}(E_w, S_{w,d,t}),
\qquad k \in \{O, I\}.
\end{align}
The baseline vector $g^{k}$ contains an intercept, the current state $E_w$, a
morning/afternoon indicator, a weekday/weekend indicator, recent burden, and
the previous decision slot's standardized page-view value (a first-order
autoregressive term); the moderator vector $f^{k}$ contains an intercept and
$E_w$, allowing the engagement effect of a walking suggestion to depend on
current perceived utility.

\paragraph{Daily mediators.}
Fitbit wear and daily check-in completion follow logistic models at the daily
level, moderated separately by the morning and afternoon suggestions: for
$k \in \{FW, PJ\}$,
\begin{align}
M_{w,d}^{E,k} \mid E_w,\, J_w^{\mathrm{week}}
\sim
\operatorname{Bernoulli}\!\bigl(\pi_{w,d}^{k}\bigr),
\qquad
\operatorname{logit}\bigl(\pi_{w,d}^{k}\bigr)
=
g^{k}(S_{w,d})^\top \alpha^{k}
+ \sum_{t=1}^{2}
A_{w,d,t}\, f_t^{k}(S_{w,d})^\top \beta_t^{k}
+
q^{k}(E_w, S_{w,d}).
\label{eq:pu-daily-models}
\end{align}
The baseline vectors contain an intercept, $E_w$, a weekday/weekend indicator,
recent burden, and the previous day's value of the same mediator (a first-order
autoregressive term); each moderator vector $f_t^{k}$ contains an intercept and
$E_w$.

\subsubsection{Estimation}
\label{sssec:pu-estimation}

Because the latent state is one-dimensional, the likelihood implied by
\eqref{eq:ssm-ew}--\eqref{eq:pu-daily-models} can be computed by deterministic
quadrature rather than by simulation. We evaluate the prediction--filtering
recursion on a fixed grid $\mathcal{G} = \{e_1, \ldots, e_K\} \subset [-3, 3]$
with trapezoidal weights: week~1 contributes its emission likelihood at the
fixed point $E_1 = 2$; each subsequent week propagates the filtered density
through the Gaussian transition \eqref{eq:ssm-ew} and multiplies in that week's
emission likelihood, accumulating the log-likelihood from the normalizing
increments. After week~$W$, the last Sunday $J_{W+1}^{\mathrm{week}}$ (and
$U_{W+1}$ if completed) is multiplied into the predictive of $E_{W+1}$; there
is no thirteenth mediator week. Each increment includes the transition mass
retained by the grid,
so probability pushed beyond the grid is penalized rather than renormalized
away; all computations are carried out in log space, and the retained and
boundary masses of the filtered density are monitored to verify that the grid
encloses the latent process.

Estimation proceeds in two stages. Stage~1 maximizes the summed penalized
log-likelihood across all participants under a single shared coefficient
vector, on a coarser grid ($K = 121$); its sole purpose is to provide a
stable, common initial value. Stage~2 fits each participant independently on a
finer grid ($K = 241$), initialized at the Stage-1 estimate. In both stages
the objective adds the ridge penalty
$\lambda \sum_k w_k (\theta_k - c_k)^2$ with $\lambda = 0.5$, where the weights
$w_k \in \{0, 1\}$ restrict the penalty to slope coefficients: all intercepts
and all log noise scales are unpenalized and are therefore estimated at their
likelihood values, so that baseline prevalences and residual variances are not
shrunk toward arbitrary reference points. In Stage~1 the center $c$ is zero.
In Stage~2 the center encodes directional working assumptions as shrinkage
targets rather than post-hoc sign constraints: $19$ components receive a nudge
of magnitude $0.05$ --- the three transition feedbacks $a_2, a_3, a_4$ ($+$),
the four $E_w$ slopes in the occurrence, intensity, wear, and check-in models
($+$), the six suggestion-by-$E_w$ interactions ($+$), and the six suggestion
main effects ($-$, reflecting that a prompt is an attentional cost absent any
offsetting utility) --- while all remaining penalized components, including the
carry-over coefficients in \eqref{eq:pu-jlag}, are shrunk toward zero.

Two devices guard the stability of the fitted latent process. First, the
persistence coefficient $a_1$ is constrained to $(-0.98, 0.98)$. Second,
because $E_w$ also feeds back on itself through the within-week mediators
($E_w \to PV_w, FW_w, PJ_w \to E_{w+1}$), we penalize a first-order proxy of
the compound loop gain,
\begin{align}
g
=
a_1
+ a_2\, \widehat{\partial PV_w / \partial E_w}
+ a_3\, \widehat{\partial FW_w / \partial E_w}
+ a_4\, \widehat{\partial PJ_w / \partial E_w},
\label{eq:pu-loop-gain}
\end{align}
where the mediator derivatives are linearizations that bound each logistic
slope by its worst case $1/4$ and combine the occurrence and intensity slopes
for the page-view hurdle. The proxy is evaluated at the four configurations
$A \in \{0, 1\} \times J_{w}^{\mathrm{week}} \in \{0, 1\}$, and the worst
case enters the smooth barrier $A_g\, g^2 / \max(1 - g^2, \varepsilon)$ with
$A_g = 3$ and $\varepsilon = 10^{-3}$, so that $\lvert g \rvert \geq 1$ ---
which would produce a saturating random walk in the simulator --- is heavily
penalized while the objective remains finite during optimization. A lighter
Fisher-$z$ barrier ($\operatorname{arctanh}^2 a_1$, weight $1$) is retained on
$a_1$ alone, so that the persistence coefficient cannot rest on its box
boundary by exploiting slack in the compound proxy. The noise scales
$\sigma_E$, $\sigma_{U_1}$, $\sigma_{U_2}$, and $\sigma_{PV}$ are parameterized
on the log scale and box-constrained to $[0.03, 10]$, with no penalty pulling
them toward any particular value.
